\section{Finite restricted masses of closed classes}
\label{sec:finite-restricted-invariants}

The preceding infinite supremum comes from the nonclosed perturbation. For a closed class, a K\"ahler modification supplies a finite comparison mass. This gives an invariant on every subvariety, including those contained in the exceptional image of a chosen modification.

Assume in this section and the next that $X$ belongs to Fujiki's class $\mathcal C$. Every irreducible reduced subvariety $Y\subset X$ has a resolution $\pi:Z\to Y$ with $Z$ compact K\"ahler. Indeed, closed subspaces of a class-$\mathcal C$ space remain in that class, and the K\"ahler modification theorem applies to $Y$. These classical geometric inputs are recorded in \cref{sec:history-closed-loci}. In particular, taking the proper transform under one fixed modification of $X$ is not needed to obtain a model of $Y$.

Two such resolutions admit a common K\"ahler refinement: resolve the closure of their common graph in $Z_1\times Z_2$ by embedded resolution. The resulting smooth strict transform is K\"ahler, since the ambient resolution is a sequence of blowups of the compact K\"ahler product along smooth centers. The K\"ahler property of a blowup follows by adding a small positive multiple of a curvature form of $\mathcal O(-E)$ to the pulled-back metric; the curvature is positive in the vertical directions along the exceptional divisor $E$.

For a closed real $(1,1)$-class $b$ on a compact K\"ahler manifold $Z$ of dimension $k>0$, set
\begin{equation}\label{eq:closed-class-mass}
 \mathcal R_Z(b)=\sup_{\substack{S\in b\\S\geq0}}
                  \int_Z\langle S^k\rangle,
 \qquad \sup\varnothing=0.
\end{equation}
For a class $a\in H^{1,1}_{\mathrm{BC}}(X,\mathbb R)$ and an irreducible reduced $k$-dimensional $Y\subset X$, write $f=j\circ\pi$, where $j:Y\hookrightarrow X$, and define
\begin{equation}\label{eq:finite-restricted-invariant}
 \mathcal R(a;Y)=\mathcal R_Z(f^*a).
\end{equation}

\begin{proposition}\label{prop:finite-restricted-invariance}
The quantities in \eqref{eq:closed-class-mass} and \eqref{eq:finite-restricted-invariant} are finite. The latter is independent of the K\"ahler resolution. If $\eta\in b$ is smooth, $\kappa$ is K\"ahler and $C>0$ satisfies $C\kappa-\eta>0$, then
\begin{equation}\label{eq:closed-class-mass-bound}
 \mathcal R_Z(b+t\{\kappa\})
       \leq(C+t)^k\int_Z\kappa^k \qquad(t\geq0).
\end{equation}
For a modification $\mu:X'\to X$ and the proper transform $Y'$ of a subvariety meeting the isomorphism locus,
\[
 \mathcal R(\mu^*a;Y')=\mathcal R(a;Y).
\]
\end{proposition}
\begin{proof}
First consider a K\"ahler form $\lambda$ on $Z$. For $q\in\PSH(Z,\lambda)$ normalized by $q\leq0$, the truncations $q_j=\max(q,-j)$ have mass $\int_Z\lambda^k$ by \cref{thm:mass-invariance-characterization}, or Stokes' theorem when $k=1$. Restricting their measures to $\{q>-j\}$ and passing to the increasing limit gives
\[
 \int_Z\langle(\lambda+\ddc q)^k\rangle
       \leq\int_Z\lambda^k.
\]
The smooth current $\lambda$ attains equality.

Put $A=C\kappa-\eta$. For a positive current $S$ in $b+t\{\kappa\}$, the current $S+A$ is positive and belongs to $(C+t)\{\kappa\}$. Positivity and multilinearity of local non-pluripolar products give $\langle S^k\rangle\leq\langle(S+A)^k\rangle$. The K\"ahler bound just proved yields \eqref{eq:closed-class-mass-bound}. The definition uses the class itself, so changing its smooth representative or the comparison form does not change the supremum.

Let $\rho:W\to Z$ be a modification between compact K\"ahler manifolds. By \cref{prop:psh-descent-modification}, every positive current in $\rho^*b$ is the pullback of a unique positive current in $b$. Their non-pluripolar measures correspond on the isomorphism locus, and both give zero mass to the exceptional analytic sets. Hence their total masses agree and
\[
 \mathcal R_W(\rho^*b)=\mathcal R_Z(b).
\]
Apply this identity on a common K\"ahler refinement to obtain independence of resolution. The same argument compares resolutions of $Y$ and $Y'$, since their maps to $X$ agree on a dense open set and therefore everywhere on the common refinement.
\end{proof}

Positivity of this mass has a classical interpretation on a K\"ahler model. We use two facts from K\"ahler pluripotential theory: a class containing a positive current with positive non-pluripolar top mass is big, and a current with minimal singularities in a big class attains its positive-product volume. For a nef class this volume is its top intersection. Their precise sources are recorded in \cref{sec:history-closed-loci}. The strict current needed for a nef class can also be obtained by the following normalized diagonal construction.

\begin{proposition}\label{prop:restricted-mass-bigness}
For a closed class $b$ on a compact K\"ahler $k$-fold $Z$,
\[
 \mathcal R_Z(b)>0\quad\Longleftrightarrow\quad b\text{ is big}.
\]
If $b$ is nef, then
\begin{equation}\label{eq:nef-restricted-mass}
 \mathcal R_Z(b)=\int_Zb^k.
\end{equation}
For nef $b$ with $I=\int_Zb^k>0$, fix a K\"ahler form $\kappa$, let $p_1,p_2:Z\times Z\to Z$ be the projections and $\Delta\subset Z\times Z$ the diagonal, and put
\[
 M=k\int_Zb^{k-1}\wedge\{\kappa\}.
\]
There is a closed positive $(k,k)$-current $\Theta$ in $(p_1^*b+p_2^*b)^k$ with $\Theta\geq d[\Delta]$, $d>0$, such that
\begin{equation}\label{eq:normalized-diagonal-current}
 K_b=M^{-1}(p_1)_*(\Theta\wedge p_2^*\kappa)\in b,
 \qquad K_b\geq(d/M)\kappa.
\end{equation}
\end{proposition}
\begin{proof}
If $\mathcal R_Z(b)>0$, some positive representative has positive top mass, so the positive-mass criterion makes $b$ big. Conversely, if $S\in b$ satisfies $S\geq\delta\kappa$, then $S-\delta\kappa$ is positive and closed. Expansion of the non-pluripolar product gives $\langle S^k\rangle\geq\delta^k\kappa^k$.

Suppose $b$ is nef. For $t>0$, the class $b+t\{\kappa\}$ is K\"ahler. Adding $t\kappa$ to an arbitrary positive representative of $b$ and using the K\"ahler mass bound yields
\[
 \mathcal R_Z(b)\leq\int_Z(b+t\{\kappa\})^k.
\]
Letting $t$ decrease to zero gives the upper bound in \eqref{eq:nef-restricted-mass}. In particular, equality holds if $I=0$.

For $I>0$, choose $L$ so that $L\{\kappa\}-b$ is K\"ahler. Nonnegativity of intersections of nef classes gives
\[
 0\leq\int_Zb^{k-1}\wedge(L\{\kappa\}-b)
      =LM/k-I,
\]
so $M>0$. The class $B=p_1^*b+p_2^*b$ is nef and satisfies
\[
 \int_{Z\times Z}B^{2k}={2k\choose k}I^2>0.
\]
The classical mass-concentration theorem along the diagonal gives $\Theta\in B^k$ with $\Theta\geq d[\Delta]$ for some $d>0$. In pushing forward $B^k\wedge p_2^*\{\kappa\}$, only the term with one factor from the first copy of $Z$ survives. Thus
\[
 (p_1)_*(B^k\wedge p_2^*\{\kappa\})=Mb,
 \qquad
 (p_1)_*([\Delta]\wedge p_2^*\kappa)=\kappa.
\]
This proves \eqref{eq:normalized-diagonal-current}. The K\"ahler $\ddc$-lemma for currents identifies the computed cohomology class with the Bott--Chern class $b$. Therefore $b$ is big, and a current with minimal singularities has mass $\int_Zb^k$. This gives the remaining lower bound.
\end{proof}

The criterion descends to a smooth subvariety $Y$. Write $a_Y=j^*a$. If $\pi:Z\to Y$ is a K\"ahler modification and $K\in\pi^*a_Y$ satisfies $K\geq\delta\kappa$, choose a Hermitian metric $h$ on $Y$ and $C>0$ with $\pi^*h\leq C\kappa$. Potential descent gives a positive representative $\pi_*K$ of $a_Y$, and
\begin{equation}\label{eq:strict-current-pushforward}
 \pi_*K\geq(\delta/C)\pi_*\pi^*h=(\delta/C)h.
\end{equation}
Conversely, a strict positive representative of $a_Y$ pulls back to a positive current whose non-pluripolar mass is positive on a coordinate ball in the isomorphism locus. On that ball it dominates a positive multiple of the closed form $\kappa$, so the local product expansion applies. Consequently, $\mathcal R(a;Y)>0$ is equivalent to bigness of the restricted class $a_Y$. This argument uses the Hermitian metric only for the current inequalities; the product comparison takes place with a closed local form.

\begin{proposition}\label{prop:restricted-rank-drop}
Let $W$ be a connected compact K\"ahler manifold of dimension $k>0$, let $g:W\to X$ be holomorphic, and suppose $\dim g(W)<k$. Then for every closed class $a$ on $X$,
\[
 \mathcal R_W(g^*a)=0.
\]
\end{proposition}
\begin{proof}
If the image is a point, $g^*a=0$ and \eqref{eq:nef-restricted-mass} gives the assertion. Otherwise, choose a K\"ahler resolution $h:Q\to g(W)$, and resolve the graph of the meromorphic lift of $g$. This gives a K\"ahler modification $\rho:\widehat W\to W$ and a holomorphic map $F:\widehat W\to Q$ with $g\circ\rho=h\circ F$, where $h$ is also viewed as a map into $X$.

Choose a smooth representative $\eta\in a$, a K\"ahler form $\kappa_Q$ on $Q$, and $C>0$ such that $C\kappa_Q-h^*\eta\geq0$. The forms
\[
 P=CF^*\kappa_Q,
 \qquad A=P-\rho^*g^*\eta
\]
are closed and semipositive. Since $\dim Q<k$, $P^k=0$. For any positive current $S\in g^*a$, the current $\rho^*S+A$ is positive in the nef class $\{P\}$. Hence
\[
 \int_{\widehat W}\langle(\rho^*S)^k\rangle
 \leq\int_{\widehat W}\langle(\rho^*S+A)^k\rangle
 \leq\mathcal R_{\widehat W}(\{P\})
 =\int_{\widehat W}P^k=0.
\]
Modification invariance proves the result. If the class has no positive representative, its invariant is zero by definition.
\end{proof}

\begin{corollary}\label{cor:finite-restricted-null-union}
Let $a$ be nef on $X$. For every positive-dimensional irreducible reduced $Y\subset X$,
\[
 \mathcal R(a;Y)=\int_Ya^{\dim Y}.
\]
If $\int_Xa^n>0$, then
\begin{equation}\label{eq:finite-restricted-null-union}
 \bigcup_{\substack{Y\subset X\ \mathrm{irreducible},\ \dim Y>0\\
                         \mathcal R(a;Y)=0}}Y
       =\operatorname{Null}(a)=E_{\mathrm{nK}}(a).
\end{equation}
\end{corollary}
\begin{proof}
Pullback preserves nefness: pull back smooth representatives with arbitrarily small negative error and bound the pulled-back Hermitian metric by a K\"ahler metric on the resolution. Now apply \eqref{eq:nef-restricted-mass}. The intersection number on $Y$ is defined by this same pullback. The first equality in \eqref{eq:finite-restricted-null-union} follows from the definition of the null locus; the second is the classical nef-and-big locus theorem recalled earlier in this chapter.
\end{proof}

When $X$ is K\"ahler and $a$ is nef, \eqref{eq:bc-numerical-restricted-volume} is also $\int_Ya^{\dim Y}$, so the two restricted invariants agree. For a general class, \eqref{eq:finite-restricted-invariant} takes the supremum over all positive currents on a resolution of $Y$; it does not require those currents to extend to $X$. The nef comparison therefore does not identify the two constructions outside that regime.
